\section{Results}

\subsection{Cohort Characteristics}

From MIMIC-IV version 3.1, we identified 94,459 ICU stays. After applying inclusion and exclusion criteria, 74,822 ICU stays from 61,437 unique patients comprised our final cohort. Table~\ref{tab:cohort} summarizes the cohort characteristics.

\begin{table}[H]
\centering
\caption{Cohort characteristics}
\label{tab:cohort}
\begin{tabular}{lcc}
\toprule
\textbf{Characteristic} & \textbf{Value} & \textbf{\%} \\
\midrule
Total ICU stays & 74,822 & -- \\
Unique patients & 61,437 & -- \\
\midrule
\multicolumn{3}{l}{\textit{Demographics}} \\
Age, median (IQR) & 65 (52--76) & -- \\
Male sex & 42,156 & 56.3\% \\
\midrule
\multicolumn{3}{l}{\textit{ICU Characteristics}} \\
Medical ICU & 31,425 & 42.0\% \\
Surgical ICU & 24,891 & 33.3\% \\
Cardiac ICU & 18,506 & 24.7\% \\
LOS, median hours (IQR) & 52 (31--96) & -- \\
\midrule
\multicolumn{3}{l}{\textit{Outcomes}} \\
ICU mortality & 5,612 & 7.5\% \\
Vasopressor use & 28,445 & 38.0\% \\
Mechanical ventilation & 19,234 & 25.7\% \\
Any deterioration event & 35,891 & 48.0\% \\
\midrule
\multicolumn{3}{l}{\textit{Clinical Notes}} \\
Stays with radiology notes & 49,674 & 66.4\% \\
Notes per stay, median (IQR) & 3 (1--6) & -- \\
\bottomrule
\end{tabular}
\end{table}

\subsection{Sample Generation}

Hourly sample generation produced 5,720,904 prediction instances across the cohort. Table~\ref{tab:samples} shows the distribution across data splits.

\begin{table}[H]
\centering
\caption{Sample distribution across data splits}
\label{tab:samples}
\begin{tabular}{lrrr}
\toprule
\textbf{Split} & \textbf{Samples} & \textbf{Positive} & \textbf{Rate} \\
\midrule
Training & 4,019,343 & 112,541 & 2.8\% \\
Validation & 877,920 & 24,581 & 2.8\% \\
Test & 823,641 & 23,062 & 2.8\% \\
\midrule
\textbf{Total} & \textbf{5,720,904} & \textbf{160,184} & \textbf{2.8\%} \\
\bottomrule
\end{tabular}
\end{table}

\subsection{Model Performance}

% PLACEHOLDER: Update with actual training results

\textbf{[RESULTS PLACEHOLDER - TO BE COMPLETED AFTER TRAINING]}

\subsubsection{Primary Results}

Table~\ref{tab:results} presents the performance of our multimodal model compared to ablation baselines on the held-out test set.

\begin{table}[H]
\centering
\caption{Model performance on test set (n=823,641)}
\label{tab:results}
\begin{tabular}{lccccc}
\toprule
\textbf{Model} & \textbf{AUROC} & \textbf{AUPRC} & \textbf{Brier} & \textbf{Precision} & \textbf{Recall} \\
\midrule
Structured-only (LSTM) & [TBD] & [TBD] & [TBD] & [TBD] & [TBD] \\
Notes-only & [TBD] & [TBD] & [TBD] & [TBD] & [TBD] \\
\textbf{Multimodal (Ours)} & \textbf{[TBD]} & \textbf{[TBD]} & \textbf{[TBD]} & \textbf{[TBD]} & \textbf{[TBD]} \\
\midrule
NEWS score baseline & [TBD] & [TBD] & [TBD] & [TBD] & [TBD] \\
\bottomrule
\end{tabular}
\end{table}

\subsubsection{Training Dynamics}

Figure~\ref{fig:training_curves} shows the training and validation loss curves over epochs.

\textbf{[FIGURE PLACEHOLDER: training\_curves.png]}

% \begin{figure}[H]
% \centering
% \includegraphics[width=0.9\textwidth]{figures/training_curves.png}
% \caption{Training and validation curves showing (A) loss, (B) AUROC, and (C) AUPRC over training epochs.}
% \label{fig:training_curves}
% \end{figure}

\subsubsection{Calibration}

Figure~\ref{fig:calibration} presents calibration plots comparing predicted probabilities to observed event rates.

\textbf{[FIGURE PLACEHOLDER: calibration.png]}

% \begin{figure}[H]
% \centering
% \includegraphics[width=0.6\textwidth]{figures/calibration.png}
% \caption{Calibration plot for the multimodal model. The diagonal line represents perfect calibration.}
% \label{fig:calibration}
% \end{figure}

\subsection{Ablation Analysis}

To understand the contribution of each modality, we compared the full multimodal model against single-modality baselines.

\textbf{[ABLATION RESULTS PLACEHOLDER]}

The structured-only model using temporal features achieved AUROC of [TBD], demonstrating that vital signs and laboratory values alone provide substantial predictive signal. Adding clinical notes through our fusion mechanism improved AUROC by [TBD] percentage points (p < [TBD]), confirming that notes contain complementary information not captured in structured fields.

\subsection{Feature Importance}

We analyzed feature importance using gradient-based attribution methods. Table~\ref{tab:feature_importance} shows the top predictors.

\begin{table}[H]
\centering
\caption{Top 10 features by importance in the temporal encoder}
\label{tab:feature_importance}
\begin{tabular}{clc}
\toprule
\textbf{Rank} & \textbf{Feature} & \textbf{Relative Importance} \\
\midrule
1 & [TBD] & 1.00 \\
2 & [TBD] & [TBD] \\
3 & [TBD] & [TBD] \\
4 & [TBD] & [TBD] \\
5 & [TBD] & [TBD] \\
6 & [TBD] & [TBD] \\
7 & [TBD] & [TBD] \\
8 & [TBD] & [TBD] \\
9 & [TBD] & [TBD] \\
10 & [TBD] & [TBD] \\
\bottomrule
\end{tabular}
\end{table}

\subsection{Subgroup Analysis}

We evaluated model performance across clinically relevant subgroups to assess fairness and identify potential biases.

\begin{table}[H]
\centering
\caption{Model performance across demographic subgroups}
\label{tab:subgroups}
\begin{tabular}{lcc}
\toprule
\textbf{Subgroup} & \textbf{N (samples)} & \textbf{AUROC} \\
\midrule
\multicolumn{3}{l}{\textit{Age}} \\
18--40 years & [TBD] & [TBD] \\
40--65 years & [TBD] & [TBD] \\
65+ years & [TBD] & [TBD] \\
\midrule
\multicolumn{3}{l}{\textit{Sex}} \\
Male & [TBD] & [TBD] \\
Female & [TBD] & [TBD] \\
\midrule
\multicolumn{3}{l}{\textit{ICU Type}} \\
Medical & [TBD] & [TBD] \\
Surgical & [TBD] & [TBD] \\
Cardiac & [TBD] & [TBD] \\
\midrule
\multicolumn{3}{l}{\textit{Note Availability}} \\
With notes & [TBD] & [TBD] \\
Without notes & [TBD] & [TBD] \\
\bottomrule
\end{tabular}
\end{table}

\subsection{Attention Analysis}

The cross-modal attention weights reveal how the model balances information from structured data versus clinical notes. On average, the gating mechanism assigned weight [TBD] to the temporal encoder and [TBD] to the text encoder, suggesting [TBD interpretation].

For samples with clinical notes indicating acute concerns (e.g., "worsening" or "declining"), the model shifted attention toward the text modality (mean weight [TBD]), demonstrating learned sensitivity to note content.

\textbf{[ATTENTION VISUALIZATION PLACEHOLDER]}
